Agent-based modeling (ABM) has developed from a specialized approach to representing heterogeneous interacting entities into a widely used framework for studying coupled social, biological, urban, and technological systems. This review combines a corpus analysis of 29,139 publications from 1972 to 2025 with established research on ABM design, large language model (LLM)-based agents, and scientific collaboration networks. Three findings emerge. First, annual publication output rose from 134 papers in 2000 to 2,495 in 2025, alongside a thematic shift from ecology and the environment toward epidemiology, public health, social dynamics, transport, and urban systems. Second, LLM adoption is measurable but remains at an early stage: the share of ABM papers using LLMs increased from 1.50\% in 2023 to 5.65\% in 2025. Of the identified uses, 59.6\% involved agent behavior or decision-making, while 40.4\% supported implementation, data generation, analysis, or interfaces. Third, collaboration became more extensive and international, yet also more modular: detected communities increased from 8 to 52 across the analyzed periods, with no evidence of general convergence. Together, these findings show how application priorities, agent architectures, and research communities are changing. They also point to a shared agenda: strengthen behavioral validation, transparency, reproducibility, and the evaluation of LLM-enabled agents.

\par\smallskip\noindent\textbf{Keywords:} agent-based modeling; generative agents; large language models; bibliometrics; coauthorship networks; community detection; complex adaptive systems